% GENERATED by build_chapters.py; edit docs/report/chapters, not this file

\section{Current System Limitations}

\subsection{What the results do not show}

\begin{itemize}[leftmargin=*, itemsep=2pt]
\item \textbf{One demo application.} All heal measurements use a small Express app with one to three source files. Bugs in a large code base, in other languages or in build configuration are not covered.
\item \textbf{Small case sets.} The quality studies use 11 to 17 cases with two or three repetitions each (Chapter 9). A difference of one case is within noise, and the report treats it that way.
\item \textbf{First-try quality is partial.} With thinking off, the production setting for the first attempt, 12 of 18 hard single-file calls were healed. Thinking raised that to 18 of 18 at about four times the time, which is why retries use it (docs/\allowbreak{}results/\allowbreak{}SHIPPED\_\allowbreak{}QUALITY\_\allowbreak{}EVAL.md).
\item \textbf{Multi-file healing is below its own bar.} The fast pass healed 9 of 18 against a bar of 12, and the retry loop reached 15 of 18 in replay. The model-comparison study recorded the same weakness as a "caller gap": bugs that need a change in a file other than the one the stack trace names (docs/\allowbreak{}results/\allowbreak{}MULTIFILE\_\allowbreak{}EVAL.md, MODEL\_\allowbreak{}BAKEOFF.md).
\item \textbf{A passing fix can still be wrong.} The gates check that the edit applies, parses and passes the repository's tests. The recorded run of Chapter 10 produced such a fix: it removed the error and dropped a string conversion.
\item \textbf{The OpenAI fallback has never produced a fix.} Its routing is tested, but on the n8n Cloud trial the free credits answered \texttt{Forbidden} to a webhook-triggered run, and the cause was not established.
\end{itemize}

\subsection{Operational limits}

\begin{itemize}[leftmargin=*, itemsep=2pt]
\item \textbf{The model address changes on every Kaggle boot.} The Set Model URL webhook removes the manual step but not the dependency: a stopped session means a dead address, and the workflow then reaches its failure comment or the fallback.
\item \textbf{GPU availability.} A session can wait in the Kaggle queue without limit; one attempt for this report waited about 90 minutes without a GPU (Chapter 10).
\item \textbf{A cap of 1.5 hours per session} in our notebook, set to bound quota use, and no way to stop a session from the command line.
\item \textbf{The n8n Cloud instance is a 14-day trial} whose REST API is disabled.
\item \textbf{The attempt cap lives in the notify job}, not in n8n. As far as the workflow's nodes show, a call to the webhook made by hand with a high \texttt{attempt} value is not stopped inside n8n.
\item \textbf{Failures that are not test failures are not healed.} A failed \texttt{npm install} or a build step that fails after the tests passed does not set \texttt{failure\_\allowbreak{}kind=code} and is ignored (Chapter 5).
\item \textbf{Seven workflow nodes have no error setting}, and 17 backend audit findings are open (Chapters 11 and 8).
\item \textbf{One healer per repository.} Running the n8n workflow and the backend on the same repository would open competing pull requests.
\end{itemize}

\subsection{Experiments that did not work}

The failures below were useful: each changed a decision.

\begingroup\small
\begin{longtable}{>{\raggedright\arraybackslash}p{0.287\dimexpr\linewidth-6\tabcolsep\relax}>{\raggedright\arraybackslash}p{0.425\dimexpr\linewidth-6\tabcolsep\relax}>{\raggedright\arraybackslash}p{0.287\dimexpr\linewidth-6\tabcolsep\relax}}
\toprule
\rowcolor{gray!14}\textbf{experiment} & \textbf{what happened} & \textbf{consequence} \\
\midrule
\endhead
\bottomrule
\endfoot
\endlastfoot
A quantization tuned for the model's DeltaNet layers (Ridge) & 5 of 18 against 9 of 18 for the generic Q4 build (docs/\allowbreak{}results/\allowbreak{}SPEED\_\allowbreak{}QUALITY\_\allowbreak{}EVAL.md) & the generic unsloth UD-Q4\_\allowbreak{}K\_\allowbreak{}XL file stayed \\
A ternary-weight model (Bonsai) & tied the Q4 build on quality at about a third of the size, but could not use tensor split or MTP & not used \\
The claim that 30 tok/\allowbreak{}s was unreachable on two T4s & refuted by measurement with tensor split and MTP (22 to 27 tok/\allowbreak{}s) & the analysis moved to docs/\allowbreak{}old/\allowbreak{} \\
Non-streamed thinking retry & died at 125.9 s on the tunnel's 120 s limit (HTTP 524) & model calls are streamed \\
Following the log redirect with the n8n HTTP node & the GitHub token would have been sent to the storage host & the redirect is handled by two nodes (Chapter 12) \\
Judging a feature by a log line & wording changes in logs twice made a disabled feature look enabled & features are judged by effect, as in the thinking self-check \\
Asking the model what it is & the answers about its own identity were unreliable & the endpoint's \texttt{/\allowbreak{}v1/\allowbreak{}models} is the reference \\
n8n 2.x & the webhook was never registered & n8n 1.x only \\
Reading \texttt{\$vars} in a self-hosted n8n & variables are a licensed feature and did not resolve & the local import substitutes the literal address \\
\bottomrule
\caption{Experiments that did not work}\label{tab:12acurren-1}\\
\end{longtable}
\endgroup

Sources for the rows: docs/\allowbreak{}research/\allowbreak{}QWEN38\_\allowbreak{}IDENTITY\_\allowbreak{}AND\_\allowbreak{}RIDGE.md, docs/\allowbreak{}results/\allowbreak{}SPEED\_\allowbreak{}QUALITY\_\allowbreak{}EVAL.md, docs/\allowbreak{}results/\allowbreak{}SPEED\_\allowbreak{}LEVERS\_\allowbreak{}EVAL.md, docs/\allowbreak{}results/\allowbreak{}GOLIVE\_\allowbreak{}RESULTS.md, tools/\allowbreak{}golive/\allowbreak{}n8n\_\allowbreak{}import.py, and tools/\allowbreak{}n8n\_\allowbreak{}rig/\allowbreak{}README.md.

\subsection{Inconsistencies found in the older documentation}

The review of the existing documents for this report found, and the rewritten documents fix, statements that no longer matched the code: the old setup guide named a different demo layout and a 20-node workflow; two documents said the backend turns thinking on for classifier re-asks when no code does; the speed analysis in \texttt{docs/\allowbreak{}old/\allowbreak{}} was refuted by the later speed study; and one endpoint note counted a crashed kernel version and its successor as two sessions. The raw data of several early experiments (the first sweep, the first end-to-end baseline) was not kept, so those numbers can be quoted but not recomputed.

\section{Future Enhancements}

The items below follow from the limits in Chapter 14, in the order we would do them.

\begin{enumerate}[leftmargin=*, itemsep=2pt]
\item \textbf{Close the open backend findings} that can lose a delivery or corrupt state: the SQLite lock held across awaits (move intake off the event loop, or use Postgres), the stuck \texttt{processing} rows, the rollback that loses a pull-request record and the unbounded rate-limit sleeps (docs/\allowbreak{}results/\allowbreak{}BACKEND\_\allowbreak{}AUDIT.md).
\item \textbf{Give every workflow node an error path.} The seven nodes without one should continue as data into the existing failure comment, so that a GitHub refusal produces a message and not a silent red run. Add the attempt cap inside n8n as well, so that it does not depend on the notify job alone.
\item \textbf{A real fallback model.} Bind the OpenAI node to a credential that works for webhook-triggered runs, or replace it with a second self-hosted endpoint, and then record a heal that the fallback actually produced. Until then the fallback is wiring, not capability.
\item \textbf{Make the endpoint key random per boot.} The notebook could generate a key at start and post it to the Set Model URL webhook together with the address, so that no secret is committed and a stale address is useless.
\item \textbf{Multi-file bugs.} The caller gap needs a retrieval step that is more than the stack trace and the changed files: for example the files that import or call the changed function. The evaluation harness already has the cases for it (eval/\allowbreak{}cases/\allowbreak{}multifile\_\allowbreak{}cases.json).
\item \textbf{Judge the fix, not only the tests.} A second model call or a simple check could compare the proposed edit with the deleted or changed lines in the breaking commit and flag a fix that removes behaviour the commit did not touch, as in the recorded run.
\item \textbf{More than one repository and language.} Run the same measurements on a second demo application and on a Python project, to see which of the numbers in Chapter 9 depend on the demo.
\item \textbf{Cover build and install failures.} Extend the notify job's classification beyond failing tests, with the same refusal for infrastructure problems.
\item \textbf{Merge the work into the main branch} of the repository once the open items above are decided. The branches \texttt{harsh/\allowbreak{}kaggle-\allowbreak{}backend-\allowbreak{}integration} and \texttt{harsh/\allowbreak{}launcher-\allowbreak{}and-\allowbreak{}demo-\allowbreak{}app} hold everything described in this report.
\end{enumerate}
